% !TeX root = main.tex
\chapter{はじまり}
今日、思いついたことを書き連ねる。

これからの目標。

一、勉強会(git)を開催する。説明することによって、自分の理解を深める。人見知りを辞める。

二、社会の課題を解決する（小さいことでも。バトミントンのやつはいいかも）

三、おりちゃんを大切にする。

四、おりちゃんとの子供について真剣に考える。



仕事についてのこと。

社会の問題を解決するために働きたい。それで対価をもらって生活をしたい。今も間接的にできているかもしれないが、実感がない。

社会はもう少しシンプルでいいと思う。

何のために技術や理論を学ぶのか。役に立てるためだ。

一つのことに熱中できないのはきちんと認める。何かにおいて一番になることを諦める。

自分の好きなこと、やりたいことを社会や仕事に押し付けない。社会や他者や自分は、何が足りないのかを見極める。

社会に対して、愚痴を言わない。しかし、足りないところを見つけるのは大事である。

話すだけでなく、実行する人間になる。多分、それが得意な気がする。
\\

社会の方向性について

生活していると、何もしなくても様々な困難に直面する。自然災害、
怪我、人間関係、病。世界は複雑なので、それに対処して生きていかないといけない。
そのために、各個人が能力を磨き、それを発揮し、協力して困難に対処する。それ以上でも以下でもない。生きるとはそれだけでいいのではないか。

互いに競争したり、憎しみ合ったり、戦う必要はない。そんな暇もないはずだ。

まずは、会社単位でそれを実現してみたい。

